\chapter{绪论}

\section{研究背景与意义}

近年来，随着移动互联网、物联网和5G通信技术的快速发展，智能手机、智能手表、智能家居等智能终端设备已深入人们的生活。据IDC统计，2025年全球智能终端设备出货量已超过50亿台。这些设备每天都在产生海量的图像、语音、文本等多模态数据，对数据进行智能化处理的需求日益增长\cite{krizhevsky2012imagenet}。

人工智能技术，特别是深度学习技术的突破性进展，使得图像分类、目标检测、自然语言处理等任务取得了前所未有的效果。然而，传统的AI应用主要依赖云端计算模式：终端设备采集数据后上传至云服务器进行推理，再将结果返回终端。这种模式虽然具有强大的计算能力，但在实际应用中面临以下突出问题：

（1）隐私安全问题。用户数据需要上传至云端进行处理，存在数据泄露和被滥用的风险。特别是在医疗影像、金融证件等敏感领域，数据传输和存储的安全性难以得到充分保障\cite{zhang2020edge}。

（2）网络延迟问题。云端推理需要数据在终端与服务器之间往返传输，网络延迟可能导致用户体验下降，尤其在自动驾驶、工业检测等实时性要求较高的场景中影响更为显著。

（3）离线可用性问题。在网络条件较差或无法联网的环境下（如地下停车场、偏远地区等），云端AI服务无法正常使用。

（4）服务器成本问题。大规模的云端AI服务需要投入大量的计算资源，企业运营成本较高。

为解决上述问题，端侧AI（On-device AI）技术应运而生。端侧AI将推理计算从云端迁移到终端设备本地执行，具有隐私保护、低延迟、离线可用和降低云端负载等优势\cite{lane2016deepx}。然而，端侧设备计算资源有限，如何在保证推理精度的前提下实现模型的高效部署，是端侧AI面临的核心技术挑战。

华为于2019年发布的鸿蒙操作系统和MindSpore深度学习框架，为端侧AI应用的开发提供了系统级解决方案。鸿蒙系统采用分布式架构设计，支持多设备协同，其ArkTS开发语言和声明式UI框架使得开发者能够高效构建跨设备应用。MindSpore作为全场景AI计算框架，提供了从模型训练到端侧推理的完整工具链，其MindSpore Lite子框架专门针对端侧设备进行了优化\cite{huawei2021harmony}。

因此，研究基于鸿蒙系统与MindSpore的端侧AI应用开发具有重要的理论意义和实用价值：一方面，探索端侧深度学习模型的轻量化部署方法，丰富端侧AI的技术体系；另一方面，为鸿蒙生态下的AI应用开发提供技术参考和实践案例。

\section{国内外研究现状}

\subsection{端侧AI技术研究现状}

端侧AI技术的研究最早可追溯到2010年代初期。早期的端侧AI研究主要集中在模型压缩和推理加速两个方面。Han等\cite{han2015deep}提出了深度压缩方法，通过剪枝、量化和霍夫曼编码将模型参数压缩至原始大小的1/35，同时保持接近原始精度。Howard等\cite{howard2017mobilenets}提出了MobileNet系列模型，利用深度可分离卷积替代标准卷积，大幅减少了计算量和参数量，使得移动端实时推理成为可能。

在推理引擎方面，Google于2017年发布了TensorFlow Lite，支持Android和iOS平台的端侧推理。Apple于2019年推出了Core ML框架和Neural Engine芯片，为iOS设备提供了高效的端侧AI能力。Facebook发布了ONNX Runtime和PyTorch Mobile，支持多种硬件平台的端侧模型部署\cite{wu2019machine}。

国内方面，华为MindSpore Lite、百度Paddle Lite、小米MACE等端侧推理框架相继推出，形成了较为完善的端侧AI工具链。其中，MindSpore Lite支持模型量化、推理优化等特性，在华为自研的NPU上具有优异的推理性能\cite{mindspore2021lite}。

\subsection{图像分类技术研究现状}

图像分类是计算机视觉领域的基础任务之一。2012年，Krizhevsky等\cite{krizhevsky2012imagenet}提出的AlexNet网络在ImageNet竞赛中取得了突破性成绩，标志着深度学习在图像分类领域的成功应用。此后，VGGNet、GoogLeNet、ResNet等经典网络相继被提出，分类精度不断提升。

He等\cite{he2016deep}提出的ResNet通过残差连接解决了深层网络的退化问题，使得训练超过100层的网络成为可能。ResNet-50作为ResNet系列中最常用的变体之一，在图像分类、特征提取等任务中被广泛使用。在轻量化模型方面，SqueezeNet、ShuffleNet、EfficientNet等网络通过精心设计网络结构，在精度和效率之间取得了较好的平衡。

\subsection{鸿蒙系统应用开发现状}

鸿蒙系统自发布以来，其应用生态快速发展。华为提供了ArkTS语言和ArkUI声明式开发框架，支持声明式、组件化的应用开发模式。在AI能力方面，鸿蒙系统集成了AI子系统，提供NAPI（Native API）接口，支持调用MindSpore Lite推理引擎进行端侧推理\cite{huawei2023arkts}。

目前，基于鸿蒙系统的AI应用开发还处于起步阶段，公开的实践案例和技术文档较少。如何在鸿蒙系统中实现深度学习模型的高效推理，如何设计友好的用户交互界面，如何充分利用鸿蒙分布式能力实现跨设备AI协同，都是需要深入研究和探索的问题。

\section{本文主要研究内容}

本文以鸿蒙系统与MindSpore为技术平台，设计并实现一个智能图像分类应用。主要研究内容包括：

\begin{enumerate}
	\item 鸿蒙系统架构分析与开发环境搭建。深入分析鸿蒙系统的分布式架构和ArkTS开发框架，搭建完整的鸿蒙应用开发环境。
	\item ResNet-50图像分类模型的训练与优化。基于MindSpore框架搭建ResNet-50模型，在CIFAR-10数据集上进行训练，并通过模型量化、知识蒸馏等策略进行轻量化优化。
	\item 鸿蒙端侧推理集成。将优化后的模型转换为MindSpore Lite格式，通过鸿蒙NAPI接口在端侧设备上实现模型推理。
	\item 应用界面设计与实现。基于鸿蒙ArkTS框架设计应用的界面布局和交互逻辑，实现图片选取、预处理、推理和结果展示的全流程。
	\item 系统测试与性能分析。在鸿蒙真机上对应用进行功能测试和性能分析，验证系统的可行性和有效性。
\end{enumerate}

\section{本文组织结构}

本文共分为六章，各章内容安排如下：

第一章为绪论，介绍研究背景与意义，综述国内外研究现状，明确本文的研究内容和组织结构。

第二章为相关技术概述，介绍鸿蒙系统架构、MindSpore框架、深度学习与图像分类的基础理论知识。

第三章为系统需求分析与总体设计，分析系统的功能需求和非功能需求，设计系统的整体架构和技术路线。

第四章为系统详细设计与实现，详细阐述模型的训练与优化、端侧推理集成、应用界面开发等核心模块的设计与实现过程。

第五章为系统测试与结果分析，介绍测试环境与方案，展示功能测试和性能测试结果，并对结果进行分析。

第六章为总结与展望，总结本文的主要工作和创新点，分析存在的不足，并对未来研究方向进行展望。